% Source: https://artofproblemsolving.com/community/c3h2035673_counting_problem
% Stable identity: aops-topic-2035673-q01
Arrangements correspond to permutations \(\sigma\) of the six positions. Write \(c(\sigma)\) for the number of cycles of \(\sigma\) (fixed points count as cycles).

\emph{Facts.} A swap changes the number of cycles by exactly \(\pm1\) (it either splits a cycle into two or merges two cycles). Starting from the identity, which has \(6\) cycles, after \(3\) swaps we have \(c(\sigma)\in\{3,5\}\) (it changes parity each time and drops by at most \(3\)). Conversely:
\begin{itemize}
\item if \(c(\sigma)=3\), then \(\sigma\) is a product of \(6-3=3\) swaps (a \(k\)-cycle is a product of \(k-1\) swaps);
\item if \(c(\sigma)=5\), \(\sigma\) is a single swap \(t\), and \(t=t\cdot u\cdot u\) for any swap \(u\), so it is reachable in exactly \(3\) switches.
\end{itemize}
So we count permutations of \(6\) objects with \(3\) or \(5\) cycles.

Permutations with \(5\) cycles are the \(\binom62=15\) swaps. Permutations with \(3\) cycles have cycle type \((4,1,1)\), \((3,2,1)\) or \((2,2,2)\):
\[
\binom64\cdot3!=90,\qquad \binom63\cdot2\cdot\binom32=120,\qquad \dfrac{6!}{2^3\cdot3!}=15,
\]
a total of \(225\). Hence the number of possible arrangements is \(225+15=\boxed{240}\).
